% Options for packages loaded elsewhere
\PassOptionsToPackage{unicode}{hyperref}
\PassOptionsToPackage{hyphens}{url}
\documentclass[
]{article}
\usepackage{xcolor}
\usepackage[margin=1in]{geometry}
\usepackage{amsmath,amssymb}
\setcounter{secnumdepth}{-\maxdimen} % remove section numbering
\usepackage{iftex}
\ifPDFTeX
  \usepackage[T1]{fontenc}
  \usepackage[utf8]{inputenc}
  \usepackage{textcomp} % provide euro and other symbols
\else % if luatex or xetex
  \usepackage{unicode-math} % this also loads fontspec
  \defaultfontfeatures{Scale=MatchLowercase}
  \defaultfontfeatures[\rmfamily]{Ligatures=TeX,Scale=1}
\fi
\usepackage{lmodern}
\ifPDFTeX\else
  % xetex/luatex font selection
\fi
% Use upquote if available, for straight quotes in verbatim environments
\IfFileExists{upquote.sty}{\usepackage{upquote}}{}
\IfFileExists{microtype.sty}{% use microtype if available
  \usepackage[]{microtype}
  \UseMicrotypeSet[protrusion]{basicmath} % disable protrusion for tt fonts
}{}
\makeatletter
\@ifundefined{KOMAClassName}{% if non-KOMA class
  \IfFileExists{parskip.sty}{%
    \usepackage{parskip}
  }{% else
    \setlength{\parindent}{0pt}
    \setlength{\parskip}{6pt plus 2pt minus 1pt}}
}{% if KOMA class
  \KOMAoptions{parskip=half}}
\makeatother
\usepackage{color}
\usepackage{fancyvrb}
\newcommand{\VerbBar}{|}
\newcommand{\VERB}{\Verb[commandchars=\\\{\}]}
\DefineVerbatimEnvironment{Highlighting}{Verbatim}{commandchars=\\\{\}}
% Add ',fontsize=\small' for more characters per line
\usepackage{framed}
\definecolor{shadecolor}{RGB}{248,248,248}
\newenvironment{Shaded}{\begin{snugshade}}{\end{snugshade}}
\newcommand{\AlertTok}[1]{\textcolor[rgb]{0.94,0.16,0.16}{#1}}
\newcommand{\AnnotationTok}[1]{\textcolor[rgb]{0.56,0.35,0.01}{\textbf{\textit{#1}}}}
\newcommand{\AttributeTok}[1]{\textcolor[rgb]{0.13,0.29,0.53}{#1}}
\newcommand{\BaseNTok}[1]{\textcolor[rgb]{0.00,0.00,0.81}{#1}}
\newcommand{\BuiltInTok}[1]{#1}
\newcommand{\CharTok}[1]{\textcolor[rgb]{0.31,0.60,0.02}{#1}}
\newcommand{\CommentTok}[1]{\textcolor[rgb]{0.56,0.35,0.01}{\textit{#1}}}
\newcommand{\CommentVarTok}[1]{\textcolor[rgb]{0.56,0.35,0.01}{\textbf{\textit{#1}}}}
\newcommand{\ConstantTok}[1]{\textcolor[rgb]{0.56,0.35,0.01}{#1}}
\newcommand{\ControlFlowTok}[1]{\textcolor[rgb]{0.13,0.29,0.53}{\textbf{#1}}}
\newcommand{\DataTypeTok}[1]{\textcolor[rgb]{0.13,0.29,0.53}{#1}}
\newcommand{\DecValTok}[1]{\textcolor[rgb]{0.00,0.00,0.81}{#1}}
\newcommand{\DocumentationTok}[1]{\textcolor[rgb]{0.56,0.35,0.01}{\textbf{\textit{#1}}}}
\newcommand{\ErrorTok}[1]{\textcolor[rgb]{0.64,0.00,0.00}{\textbf{#1}}}
\newcommand{\ExtensionTok}[1]{#1}
\newcommand{\FloatTok}[1]{\textcolor[rgb]{0.00,0.00,0.81}{#1}}
\newcommand{\FunctionTok}[1]{\textcolor[rgb]{0.13,0.29,0.53}{\textbf{#1}}}
\newcommand{\ImportTok}[1]{#1}
\newcommand{\InformationTok}[1]{\textcolor[rgb]{0.56,0.35,0.01}{\textbf{\textit{#1}}}}
\newcommand{\KeywordTok}[1]{\textcolor[rgb]{0.13,0.29,0.53}{\textbf{#1}}}
\newcommand{\NormalTok}[1]{#1}
\newcommand{\OperatorTok}[1]{\textcolor[rgb]{0.81,0.36,0.00}{\textbf{#1}}}
\newcommand{\OtherTok}[1]{\textcolor[rgb]{0.56,0.35,0.01}{#1}}
\newcommand{\PreprocessorTok}[1]{\textcolor[rgb]{0.56,0.35,0.01}{\textit{#1}}}
\newcommand{\RegionMarkerTok}[1]{#1}
\newcommand{\SpecialCharTok}[1]{\textcolor[rgb]{0.81,0.36,0.00}{\textbf{#1}}}
\newcommand{\SpecialStringTok}[1]{\textcolor[rgb]{0.31,0.60,0.02}{#1}}
\newcommand{\StringTok}[1]{\textcolor[rgb]{0.31,0.60,0.02}{#1}}
\newcommand{\VariableTok}[1]{\textcolor[rgb]{0.00,0.00,0.00}{#1}}
\newcommand{\VerbatimStringTok}[1]{\textcolor[rgb]{0.31,0.60,0.02}{#1}}
\newcommand{\WarningTok}[1]{\textcolor[rgb]{0.56,0.35,0.01}{\textbf{\textit{#1}}}}
\usepackage{graphicx}
\makeatletter
\newsavebox\pandoc@box
\newcommand*\pandocbounded[1]{% scales image to fit in text height/width
  \sbox\pandoc@box{#1}%
  \Gscale@div\@tempa{\textheight}{\dimexpr\ht\pandoc@box+\dp\pandoc@box\relax}%
  \Gscale@div\@tempb{\linewidth}{\wd\pandoc@box}%
  \ifdim\@tempb\p@<\@tempa\p@\let\@tempa\@tempb\fi% select the smaller of both
  \ifdim\@tempa\p@<\p@\scalebox{\@tempa}{\usebox\pandoc@box}%
  \else\usebox{\pandoc@box}%
  \fi%
}
% Set default figure placement to htbp
\def\fps@figure{htbp}
\makeatother
\setlength{\emergencystretch}{3em} % prevent overfull lines
\providecommand{\tightlist}{%
  \setlength{\itemsep}{0pt}\setlength{\parskip}{0pt}}
\usepackage{bookmark}
\IfFileExists{xurl.sty}{\usepackage{xurl}}{} % add URL line breaks if available
\urlstyle{same}
\hypersetup{
  pdftitle={G-methods{,} IV{,} DiD{,} RDD; heterogeneity of treatment effect},
  hidelinks,
  pdfcreator={LaTeX via pandoc}}

\title{G-methods, IV, DiD, RDD; heterogeneity of treatment effect}
\author{}
\date{\vspace{-2.5em}}

\begin{document}
\maketitle

\textbf{Workflow lab:} Goal → Approach → Execution → Check → Report.

\subsection{Learning objectives}\label{learning-objectives}

\begin{itemize}
\tightlist
\item
  Use inverse-probability weighting as a g-method for time-varying
  confounding.
\item
  Recognise the data-generating situations that justify an
  instrumental-variable (IV), difference-in-differences (DiD), or
  regression-discontinuity (RDD) design.
\item
  Estimate and interpret heterogeneity of treatment effect (HTE) with a
  subgroup-agnostic method.
\end{itemize}

\subsection{Prerequisites}\label{prerequisites}

Propensity-score and IPTW basics (W3 S4).

\subsection{Background}\label{background}

Traditional regression adjustment fails for time-varying confounders
that are affected by prior treatment --- adjusting for them blocks part
of the treatment effect, failing to adjust leaves residual confounding.
Robins's g-methods solve this by reweighting or standardising in a way
that respects the temporal order. IV methods exploit a variable that
affects treatment but not outcome directly; DiD compares changes over
time between groups sharing a counterfactual trend; RDD exploits a sharp
assignment rule around a threshold.

Heterogeneity of treatment effect is the shift from ``is the average
effect non-zero?'' to ``for whom is it largest?'' Model-based HTE uses
causal forests, meta-learners (S-, T-, X-learner), or Bayesian
hierarchical models to estimate conditional average treatment effects
(CATEs) while controlling for multiple testing.

\subsection{Setup}\label{setup}

\begin{Shaded}
\begin{Highlighting}[]
\FunctionTok{library}\NormalTok{(tidyverse)}
\FunctionTok{library}\NormalTok{(broom)}
\FunctionTok{set.seed}\NormalTok{(}\DecValTok{42}\NormalTok{)}
\FunctionTok{theme\_set}\NormalTok{(}\FunctionTok{theme\_minimal}\NormalTok{(}\AttributeTok{base\_size =} \DecValTok{12}\NormalTok{))}
\end{Highlighting}
\end{Shaded}

\subsection{1. Goal}\label{goal}

Simulate a time-varying-confounding scenario, estimate the causal effect
naively and then with IPTW (a marginal structural model), and finish
with a simple HTE estimate.

\subsection{2. Approach}\label{approach}

\begin{Shaded}
\begin{Highlighting}[]
\NormalTok{df }\OtherTok{\textless{}{-}} \FunctionTok{tibble}\NormalTok{(}
  \AttributeTok{l0 =} \FunctionTok{rnorm}\NormalTok{(n),}
  \AttributeTok{a0 =} \FunctionTok{rbinom}\NormalTok{(n, }\DecValTok{1}\NormalTok{, }\FunctionTok{plogis}\NormalTok{(}\FloatTok{0.2} \SpecialCharTok{*}\NormalTok{ l0)),}
  \AttributeTok{l1 =} \FunctionTok{rnorm}\NormalTok{(n, }\AttributeTok{mean =} \FloatTok{0.5} \SpecialCharTok{*}\NormalTok{ a0 }\SpecialCharTok{+} \FloatTok{0.3} \SpecialCharTok{*}\NormalTok{ l0),}
  \AttributeTok{a1 =} \FunctionTok{rbinom}\NormalTok{(n, }\DecValTok{1}\NormalTok{, }\FunctionTok{plogis}\NormalTok{(}\FloatTok{0.3} \SpecialCharTok{*}\NormalTok{ l1 }\SpecialCharTok{+} \FloatTok{0.2} \SpecialCharTok{*}\NormalTok{ a0)),}
  \AttributeTok{y  =} \FloatTok{1.0} \SpecialCharTok{*}\NormalTok{ a0 }\SpecialCharTok{+} \FloatTok{0.8} \SpecialCharTok{*}\NormalTok{ a1 }\SpecialCharTok{+} \FloatTok{0.5} \SpecialCharTok{*}\NormalTok{ l0 }\SpecialCharTok{+} \FloatTok{0.5} \SpecialCharTok{*}\NormalTok{ l1 }\SpecialCharTok{+} \FunctionTok{rnorm}\NormalTok{(n)}
\NormalTok{)}
\end{Highlighting}
\end{Shaded}

The true effect of each treatment is positive and additive. \texttt{l1}
is a time-varying confounder because it is affected by \texttt{a0} and
predicts both \texttt{a1} and \texttt{y}.

\subsection{3. Execution}\label{execution}

Naive regression adjusting for \texttt{l1} blocks part of the
\texttt{a0} effect.

\begin{Shaded}
\begin{Highlighting}[]
\FunctionTok{tidy}\NormalTok{(naive, }\AttributeTok{conf.int =} \ConstantTok{TRUE}\NormalTok{) }\SpecialCharTok{|\textgreater{}}
  \FunctionTok{filter}\NormalTok{(term }\SpecialCharTok{\%in\%} \FunctionTok{c}\NormalTok{(}\StringTok{"a0"}\NormalTok{, }\StringTok{"a1"}\NormalTok{))}
\end{Highlighting}
\end{Shaded}

IPTW with stabilised weights:

\begin{Shaded}
\begin{Highlighting}[]
\NormalTok{den0 }\OtherTok{\textless{}{-}} \FunctionTok{glm}\NormalTok{(a0 }\SpecialCharTok{\textasciitilde{}}\NormalTok{ l0, }\AttributeTok{data =}\NormalTok{ df, }\AttributeTok{family =}\NormalTok{ binomial)}
\NormalTok{num1 }\OtherTok{\textless{}{-}} \FunctionTok{glm}\NormalTok{(a1 }\SpecialCharTok{\textasciitilde{}}\NormalTok{ a0, }\AttributeTok{data =}\NormalTok{ df, }\AttributeTok{family =}\NormalTok{ binomial)}
\NormalTok{den1 }\OtherTok{\textless{}{-}} \FunctionTok{glm}\NormalTok{(a1 }\SpecialCharTok{\textasciitilde{}}\NormalTok{ a0 }\SpecialCharTok{+}\NormalTok{ l0 }\SpecialCharTok{+}\NormalTok{ l1, }\AttributeTok{data =}\NormalTok{ df, }\AttributeTok{family =}\NormalTok{ binomial)}

\NormalTok{w0 }\OtherTok{\textless{}{-}} \FunctionTok{ifelse}\NormalTok{(df}\SpecialCharTok{$}\NormalTok{a0 }\SpecialCharTok{==} \DecValTok{1}\NormalTok{, }\FunctionTok{fitted}\NormalTok{(num0), }\DecValTok{1} \SpecialCharTok{{-}} \FunctionTok{fitted}\NormalTok{(num0)) }\SpecialCharTok{/}
      \FunctionTok{ifelse}\NormalTok{(df}\SpecialCharTok{$}\NormalTok{a0 }\SpecialCharTok{==} \DecValTok{1}\NormalTok{, }\FunctionTok{fitted}\NormalTok{(den0), }\DecValTok{1} \SpecialCharTok{{-}} \FunctionTok{fitted}\NormalTok{(den0))}
\NormalTok{w1 }\OtherTok{\textless{}{-}} \FunctionTok{ifelse}\NormalTok{(df}\SpecialCharTok{$}\NormalTok{a1 }\SpecialCharTok{==} \DecValTok{1}\NormalTok{, }\FunctionTok{fitted}\NormalTok{(num1), }\DecValTok{1} \SpecialCharTok{{-}} \FunctionTok{fitted}\NormalTok{(num1)) }\SpecialCharTok{/}
      \FunctionTok{ifelse}\NormalTok{(df}\SpecialCharTok{$}\NormalTok{a1 }\SpecialCharTok{==} \DecValTok{1}\NormalTok{, }\FunctionTok{fitted}\NormalTok{(den1), }\DecValTok{1} \SpecialCharTok{{-}} \FunctionTok{fitted}\NormalTok{(den1))}
\NormalTok{df}\SpecialCharTok{$}\NormalTok{sw }\OtherTok{\textless{}{-}}\NormalTok{ w0 }\SpecialCharTok{*}\NormalTok{ w1}

\NormalTok{msm }\OtherTok{\textless{}{-}} \FunctionTok{lm}\NormalTok{(y }\SpecialCharTok{\textasciitilde{}}\NormalTok{ a0 }\SpecialCharTok{+}\NormalTok{ a1, }\AttributeTok{data =}\NormalTok{ df, }\AttributeTok{weights =}\NormalTok{ sw)}
\FunctionTok{tidy}\NormalTok{(msm, }\AttributeTok{conf.int =} \ConstantTok{TRUE}\NormalTok{) }\SpecialCharTok{|\textgreater{}}
  \FunctionTok{filter}\NormalTok{(term }\SpecialCharTok{\%in\%} \FunctionTok{c}\NormalTok{(}\StringTok{"a0"}\NormalTok{, }\StringTok{"a1"}\NormalTok{))}
\end{Highlighting}
\end{Shaded}

\subsection{4. Check}\label{check}

Balance and weight sanity:

\begin{Shaded}
\begin{Highlighting}[]

\end{Highlighting}
\end{Shaded}

Stabilised weights should centre on 1 with no extreme tail. If they do
not, a term in the treatment model is likely misspecified.

\subsection{5. Report}\label{report}

\begin{quote}
In a simulated cohort with time-varying confounding (\emph{n} = 2000),
naive regression underestimated the effect of initial treatment. An
inverse-probability-weighted marginal structural model recovered effects
close to the true data-generating values. Instrumental variable,
difference-in-differences, and regression-discontinuity designs exploit
alternative identification strategies, each with its own assumptions.
Heterogeneity of treatment effect can be explored with causal forests or
meta-learners.
\end{quote}

\subsubsection{Brief sketches}\label{brief-sketches}

\begin{itemize}
\tightlist
\item
  \textbf{IV:} regress \texttt{y} on an instrument \texttt{z} to recover
  \texttt{β\ =\ cov(y,z)/cov(a,z)}.
\item
  \textbf{DiD:} fit \texttt{lm(y\ \textasciitilde{}\ time*group)} on
  repeated measures.
\item
  \textbf{RDD:} fit local linear regression on each side of the
  assignment cutoff.
\end{itemize}

\subsection{Common pitfalls}\label{common-pitfalls}

\begin{itemize}
\tightlist
\item
  Using ordinary regression adjustment for time-varying confounders.
\item
  Trimming extreme weights silently, hiding model misspecification.
\item
  Reporting subgroup p-values without multiplicity control as HTE
  evidence.
\end{itemize}

\subsection{Further reading}\label{further-reading}

\begin{itemize}
\tightlist
\item
  Hernán MA, Robins JM. \emph{Causal Inference: What If.}
\item
  Athey S, Imbens GW. (2019). \emph{Machine Learning Methods That
  Economists Should Know About.}
\end{itemize}

\subsection{Session info}\label{session-info}

\begin{Shaded}
\begin{Highlighting}[]

\end{Highlighting}
\end{Shaded}

\subsection{See also --- chapter index}\label{see-also-chapter-index}

\begin{itemize}
\tightlist
\item
  \href{../index.Rmd}{Methods in this chapter}
\item
  \href{../../../ch00-find-your-method.Rmd}{Find your method (Chapter
  0)}
\end{itemize}

\end{document}
